\section*{Chapter \ref{ch:b2:batteries-electrolysis} --- Batteries, Fuel Cells and Electrolysis}

\begin{solution}{exo:b2:batteries-electrolysis:1}
Zinc alone: about \qty{-0.82}{V} and a current of 0.15 (model units). Touching
copper: about \qty{-0.74}{V} and 2.4. The current, hence the dissolution rate, is
some sixteen times larger in contact with copper.
\end{solution}

\begin{solution}{exo:b2:batteries-electrolysis:2}
$m = 500 \times 3600 \times 65.38/(2 \times 96\,485) = \qty{610}{g}$.
\end{solution}

\begin{solution}{exo:b2:batteries-electrolysis:3}
$E^\circ = \qty{2.04}{V}$ (as in the chapter); six cells in series give
\qty{12}{V}.
\end{solution}

\begin{solution}{exo:b2:batteries-electrolysis:4}
$1/6.94 = \qty{0.144}{mol}$; $0.144 \times 96\,485 = \qty{1.39e4}{C}$, that is
\qty{3.86}{A.h}.
\end{solution}

\begin{solution}{exo:b2:batteries-electrolysis:5}
At 100~\%: $200 \times 86\,400 \times 63.546/(2 \times 96\,485) = \qty{5.69}{kg}$;
$\eta_F = 5.50/5.69 = 0.967$.
\end{solution}

\begin{solution}{exo:b2:batteries-electrolysis:6}
$w = 2 \times 96\,485 \times 3.3/(0.90 \times 0.06538) = \qty{1.08e7}{J/kg}$, that is
\qty{3.0}{kW.h/kg}.
\end{solution}

\begin{solution}{exo:b2:batteries-electrolysis:7}
$U = E - rI$: $r = (1.55 - 1.40)/(0.60 - 0.10) = \qty{0.30}{\ohm}$; $E = 1.55 + 0.30
\times 0.10 = \qty{1.58}{V}$.
\end{solution}

\begin{solution}{exo:b2:batteries-electrolysis:8}
Anode \ce{2Cl- -> Cl2 + 2e-}; cathode \ce{2H2O + 2e- -> H2 + 2OH-}. $\Delta_r G^\circ =
2(-157.24) - 2(-131.23) - 2(-237.13) = \qty{422.2}{kJ/mol}$, $U_{\min} =
\qty{2.19}{V}$. Chlorine would react with the hydroxide (\ce{Cl2 + 2OH- -> Cl- +
ClO- + H2O}), spoiling both products, and a mixture of chlorine and hydrogen
can explode.
\end{solution}

\begin{solution}{exo:b2:batteries-electrolysis:9}
\ce{NaCl -> Na + 1/2Cl2}, $n = 1$: $U_{\min} = 310\,674/96\,485 = \qty{3.22}{V}$.
Per kilogram of sodium: $96\,485 \times 3.22/0.02299 = \qty{1.35e7}{J}$, that is
\qty{3.75}{kW.h}.
\end{solution}

\begin{solution}{exo:b2:batteries-electrolysis:10}
On a zinc cathode the reduction of \ce{H+} is slow and its wall lies below the
zinc wave: as the potential is lowered, zinc is deposited first, with little
hydrogen. On platinum the reduction of \ce{H+} is fast and starts near
\qty{0}{V}: at first only hydrogen is evolved. Once the platinum is covered with
zinc, the cathode behaves as a zinc cathode and zinc deposits.
\end{solution}

\begin{solution}{exo:b2:batteries-electrolysis:11}
At \qty{1250}{K}: $\Delta_f G^\circ(\ce{Al2O3}) = -1328.3 + 0.75 \times 66.0 =
\qty{-1278.8}{kJ/mol}$, $\Delta_f G^\circ(\ce{CO2}) = \qty{-396.1}{kJ/mol}$. $\Delta_r G^\circ =
3(-396.1) - 2(-1278.8) = \qty{1369}{kJ}$, $n = 12$: $U_{\min} = \qty{1.18}{V}$. Inert anode:
$\Delta_r G^\circ = \qty{2557.5}{kJ}$, $U_{\min} = \qty{2.21}{V}$; the burning carbon
supplies about \qty{1.03}{V} of the work.
\end{solution}

\begin{solution}{exo:b2:batteries-electrolysis:12}
Per mole of water the charge $2F$ is the same both ways: efficiency $0.70/1.80 =
0.39$. The rest is heat: the \omterm{def:b2:current-potential-curves:overpotential}{overpotentials} of both oxygen electrodes (slow in
both directions), those of the hydrogen electrodes, and the ohmic losses.
\end{solution}

\begin{solution}{pb:b2:batteries-electrolysis:1}
\textbf{1.} Aluminium $+3 \to 0$; oxygen stays $-2$; carbon $0 \to +4$.
\textbf{2.} Cathode \ce{Al^3+ + 3e- -> Al}; anode \ce{C + 2O^2- -> CO2 + 4e-}.
\textbf{3.} \ce{2Al2O3 + 3C -> 4Al + 3CO2}, $n = 12$.
\textbf{4.} \ce{Al2O3}: $-1328.3 + 0.75 \times (1328.3 - 1262.3) = \qty{-1278.8}{kJ/mol}$;
\ce{CO2}: \qty{-396.1}{kJ/mol}.
\textbf{5.} $3(-396.1) - 2(-1278.8) = \qty{+1369}{kJ}$.
\textbf{6.} $1\,369\,000/(12 \times 96\,485) = \qty{1.18}{V}$.
\textbf{7.} $2 \times 1278.8/(12 \times 96.485) = \qty{2.21}{V}$: the oxidation of carbon,
itself \omterm{def:b2:reaction-free-energy:exergonic}{exergonic}, pays nearly half of the work of decomposing the oxide.
\textbf{8.} $10^6/26.982 = \qty{3.706e4}{mol}$.
\textbf{9.} $3 \times 96\,485 \times 3.706 \times 10^4 = \qty{1.073e10}{C}$.
\textbf{10.} $1.073 \times 10^{10}/0.94 = \qty{1.141e10}{C}$.
\textbf{11.} $1.141 \times 10^{10}/3.0 \times 10^5 = \qty{3.80e4}{s}$, \qty{10.6}{h}.
\textbf{12.} $24/10.6 = 2.27$: \qty{2.27}{t} per day.
\textbf{13.} Part of the aluminium dissolved in the bath reaches the anode gas and
is reoxidised by carbon dioxide, so its charge is spent twice; some current
also leaks through the bath without electrolysing it.
\textbf{14.} $1.141 \times 10^{10} \times 4.1 = \qty{4.68e10}{J}$, \qty{13.0}{MW.h}.
\textbf{15.} \qty{13.0}{kW.h/kg}.
\textbf{16.} $3 \times 96\,485 \times 1.18/0.026982 = \qty{1.27e7}{J/kg}$, \qty{3.5}{kW.h/kg}.
\textbf{17.} Into heat: the \omterm{def:b2:current-potential-curves:overpotential}{overpotentials} of the electrodes and, mostly, the
ohmic drop in the bath. That heat keeps the bath molten at about
\qty{1250}{K}, which would otherwise need fuel.
\textbf{18.} $4.1 \times 3.0 \times 10^5 = \qty{1.23}{MW}$.
\textbf{19.} $28\,690/0.026982 = \qty{1.06e6}{J/kg}$, \qty{0.30}{kW.h/kg}: about 2~\% of
the energy of the \omterm{def:b2:batteries-electrolysis:electrolysis}{electrolysis}.
\textbf{20.} $1.18/4.1 = 0.29$.
\textbf{21.} $3.706 \times 10^4 \times 3/4 \times 12.011 = \qty{334}{kg}$.
\textbf{22.} $3.706 \times 10^4 \times 3/4 \times 44.009 = \qty{1.22e3}{kg}$.
\textbf{23.} The hot tops of the anodes burn in the air; part of the carbon leaves
as carbon monoxide (\ce{C + CO2 -> 2CO} at that temperature), which takes twice as
much carbon per oxygen atom.
\textbf{24.} \qty{1.22}{kg} of \ce{CO2} per kilogram of aluminium.
\textbf{25.} The minimum voltage would rise by about \qty{1.03}{V}, but the anode
would give off oxygen instead of carbon dioxide and would not be consumed.
\textbf{26.} The cell uses $\boldsymbol{\approx \qty{13}{kW.h}}$ of electrical
energy per kilogram of aluminium.
\end{solution}
